\documentclass[a4paper,12pt]{article}
\usepackage[spanish,activeacute]{babel}
\usepackage[latin1]{inputenc}

%% Para las columnas
\usepackage{multicol}
\usepackage{amssymb,amsmath,eepic,fancyhdr}
\usepackage{latexsym}
\usepackage{datenumber}


%% Para numeraci?n autom?tica
\newcounter{nroproblema}
\setcounter{nroproblema}{1}
\newcounter{nroapartado}
\setcounter{nroapartado}{1}

\newcommand{\n}{%
    \vspace{.1in}\noindent{\bf \arabic{nroproblema}.\ }%
    \addtocounter{nroproblema}{1}%
    \setcounter{nroapartado}{1}%
    }%

\newcommand{\nn}{%
    \vspace{.1in}\noindent \hskip.4cm {\bf \arabic{nroproblema}.\ }%
    \addtocounter{nroproblema}{1}%
    \setcounter{nroapartado}{1}%
    }%


\newcommand{\ap}{\vspace*{0.2cm}\hspace*{0.2cm}%
   {\bf \alph{nroapartado})}
   \addtocounter{nroapartado}{1}%
   }

 
%marcos by Marco
\newcommand{\pd}[1]{\frac{\partial}{\partial #1}} %\pd{x}
\newcommand{\NN}{\ensuremath{\mathbb N}}
\newcommand{\ZZ}{\ensuremath{\mathbb Z}}
\newcommand{\CC}{\ensuremath{\mathbb C}}
\newcommand{\RR}{\ensuremath{\mathbb R}}
\newcommand{\TT}{\ensuremath{\mathbb T}}
\newcommand{\g}{\ensuremath{\frak{g}}}
\newcommand{\h}{\ensuremath{\frak{h}}}
\newcommand{\ft}{\ensuremath{\frak{t}}}
\newcommand{\vX}{\frak{X}}
\newcommand{\cB}{\mathcal{B}}
\newcommand{\cN}{\mathcal{N}}
\newcommand{\cL}{\mathcal{L}}
\newcommand{\cD}{\mathcal{D}} 
 

%% m�rgenes

% Anterior:

%\setlength{\unitlength}{1mm} \addtolength{\voffset}{-15mm}
%\addtolength{\textheight}{30mm} \addtolength{\hoffset}{-20mm}
%\addtolength{\textwidth}{45mm}

\setlength{\unitlength}{1mm} \addtolength{\voffset}{-22mm}
\addtolength{\textheight}{40mm} \addtolength{\hoffset}{-22mm}
\addtolength{\textwidth}{46mm}




\begin{document}

\noindent
{\sffamily\large Marco Zambon 
\hfill  Master UAM, 2013-14  \\Geometr\'ia simpl\'ectica y de Poisson}


 

\begin{center}\bfseries\large\textsf{Hoja 4    \vspace{-2mm}}
\end{center}
 
 
\n Considera el espacio vectorial simpl\'ectico  $(\RR^{2n},\omega )$, donde $\omega$ es la forma simpl\'ectica lineal can\'onica..

\ap Describe explicitamente el \'algebra de Lie de $Sp(2n):=\{A\in GL(2n,\RR):A^*\omega =\omega \}$.

\ap Calcula la dimensi\'on de $Sp(2n)$ (de manera distinta de la utilizada en clase).


 

%\n
%Demuestra que la siguiente es una biyecci\'on: $$\Phi \colon S^3\to SU(2), (\alpha,\beta) \mapsto \left(\begin{array}{cc} \alpha & \beta \\ -\bar{\beta} & \bar{\alpha} \end{array}\right).$$
%Aqu\'i consideramos $S^3$ como el conjunto de vectores de norma uno en $\CC^2$.

%\textbf{Observaci\'on:}  Esta aplicaci\'on es un isomorfismo de grupos de Lie. En particular, $SU(2)$ es simplemente con\'exo. 

\n [3 puntos]
Sean $\g$ una algebra de Lie y $\h$ un ideal de Lie, es decir, un subespacio tal que $[\h,\g]\subset \h$. Demuestra que el espacio cociente $\g/\h$ admite una \'unica estructura de \'algebra de Lie tal que la  projeci\'on 
$\pi \colon \g \to \g/\h$  es un morfismo de \'algebras de Lie.    



\n [3 puntos]
Sean $G,H$ grupos de Lie y  $\Phi \colon G \to H$ un  homomorfismo de grupos de Lie. Demuestra que para cada elemento $v$  del \'algebra de Lie $T_eG$ se cumple 
$$\Phi(exp(v))=exp(d_e\Phi (v))$$
donde $exp$ denota la aplicaci\'on exponential de $G$ o de $H$.
 
 
\n   Considera el \'algebra de Lie $\g:=\mathfrak{sl}_2(\RR)$, y una base
\[  \begin{array}{ccc}
h=\left( \begin{array}{cc} 1 & 0 \\ 0 & -1 \end{array} \right), &
e=\left( \begin{array}{cc} 0 & 1 \\ 0 & 0 \end{array} \right), &
f=\left( \begin{array}{cc} 0 & 0 \\ 1 & 0 \end{array} \right).
\end{array} \]


\ap Combrueba que $[h,e]=2e, [h,f]=-2f, [e,f]=h$.

\ap Calcula $d_{\g}\colon \g^*\to \wedge^2\g^*$

\ap Calcula $H^1(\g)$ y $H^2(\g)$.

\n Sea $G$ in grupo de Lie, $\g$ su \'algebra de Lie. Sea $B\colon \g\times \g \to \RR$ un producto escalar invariante bajo $Ad$, es decir, $B(Ad_gX,Ad_gY)=B(X,Y)$ para todo $g\in G$, $X,Y\in \g$.

Demuestra que, bajo el isomorfismo $\g\to \g^*, X\to B(X,\cdot)$,
la acci\'on adjunta $Ad$ de $G$ en $\g$ y la 
la acci\'on coadjunta $Ad^*$ de $G$ en $\g^*$ coinciden.
 
     \end{document}
   
  \noindent\textbf{Consejo:}
\noindent\textbf{Nota:}